\par\medskip
\hypertarget{native-073-article055}{}\pdfbookmark[2]{55.}{article-055}
\begin{center}55.\end{center}
Designando ut supra longitudinem nodi ascendentis orbitae in ecliptica per $\AscendingNode$, inclinationem per $i$; porro longitudinem nodi ascendentis plani novi in ecliptica per $n$, inclinationem per $\epsilon$; distantiam nodi ascendentis orbitae in plano novo a nodo ascendente plani novi in ecliptica (arcum $n\AscendingNode'$ in Fig. 2) per $\AscendingNode'$, inclinationem orbitae ad planum novum per $i'$; denique arcum ab $\AscendingNode$ ad $\AscendingNode'$ secundum directionem motus per $\Delta$: erunt trianguli sphaerici nostri latera $\AscendingNode-n$, $\AscendingNode'$, $\Delta$, angulique oppositi $i'$, $180^0-i$, $\epsilon$. Hinc erit secundum formulas art. praec.
\SourceLeaf{074}
\[
\sin\tfrac{1}{2}i'\sin\tfrac{1}{2}(\AscendingNode'+\Delta) = \sin\tfrac{1}{2}(\AscendingNode-n)\sin\tfrac{1}{2}(i+\epsilon)
\]
\[
\sin\tfrac{1}{2}i'\cos\tfrac{1}{2}(\AscendingNode'+\Delta) = \cos\tfrac{1}{2}(\AscendingNode-n)\sin\tfrac{1}{2}(i-\epsilon)
\]
\[
\cos\tfrac{1}{2}i'\sin\tfrac{1}{2}(\AscendingNode'-\Delta) = \sin\tfrac{1}{2}(\AscendingNode-n)\cos\tfrac{1}{2}(i+\epsilon)
\]
\[
\cos\tfrac{1}{2}i'\cos\tfrac{1}{2}(\AscendingNode'-\Delta) = \cos\tfrac{1}{2}(\AscendingNode-n)\cos\tfrac{1}{2}(i-\epsilon).
\]
Duae primae aequationes suppeditabunt $\frac{1}{2}(\AscendingNode'+\Delta)$ atque
$\sin\frac{1}{2}i'$; duae reliquae $\frac{1}{2}(\AscendingNode'-\Delta)$ atque $\cos\frac{1}{2}i'$;
ex $\frac{1}{2}(\AscendingNode'+\Delta)$ et $\frac{1}{2}(\AscendingNode'-\Delta)$ demanabunt $\AscendingNode'$ et
$\Delta$; ex $\sin\frac{1}{2}i'$ aut $\cos\frac{1}{2}i'$ (quorum consensus calculo
confirmando inserviet) prodibit $i'$. Ambiguitas, utrum $\frac{1}{2}(\AscendingNode'+\Delta)$ et
$\frac{1}{2}(\AscendingNode'-\Delta)$ inter $0$ et $180^0$ vel inter $180^0$ et $360^0$ accipere
oporteat, ita tolletur, ut tum $\sin\frac{1}{2}i'$ tum $\cos\frac{1}{2}i'$ fiant positivi,
quoniam per rei naturam $i'$ infra $180^0$ cadere debet.
